\documentclass{article}
\usepackage[utf8]{inputenc}
\usepackage[T1]{fontenc}
\usepackage{graphicx}
\usepackage{hyperref}
\usepackage{geometry}
\geometry{margin=1in}

\title{MECHANICAL ENGINEERING and CRANES (Plate CCX)}
\date{Encyclopaedia Britannica (1771)}

\begin{document}
\maketitle

\section*{Historical source}
\textbf{Entry ID:} \texttt{entry-britannica-volume-149983206-row-2077-2077} \\
\textbf{Volume ID:} 149983206 \\
\textbf{Character Count:} 867

\subsection*{Extracted Text}
[HISTORICAL COPPERPLATE ENGRAVING: MECHANICAL ENGINEERING and CRANES (Plate CCX)]
Category: Mechanics and Civil Engineering
Visual Description: Copperplate technical drawing of harbor cranes, compound pulley systems, capstans, and gear trains for heavy lifting and maritime cargo handling.

[Nemotron Super VL Analysis]: The provided image appears to be a blank sheet of aged paper with no visible engraving, text, headings, or figures. Therefore, it is not possible to transcribe any headings or describe key figures related to Plate CCX or cranes.
Identified Engraved Figures:
  * Fig. 1: Harbor crane with treadwheel and reduction gear train
  * Fig. 2: Compound block-and-tackle pulley arrangement
  * Fig. 3: Ratchet and pawl mechanism for backstop safety
Related Scientific Subjects: Crane (machine), Block and tackle, Mechanical engineering, Capstan (nautical)

\subsection*{Matched Wikipedia Article}
\textbf{Title:} History of the Encyclopædia Britannica \\
\textbf{Page ID:} 9004976 \\
\textbf{Revision ID:} 1340071003 \\
\textbf{URL:} \url{https://en.wikipedia.org/wiki/History\_of\_the\_Encyclop\%C3\%A6dia\_Britannica}


\section*{Source record}
\begin{itemize}
\item \textbf{Dataset:} \texttt{NationalLibraryOfScotland/encyclopaedia\_britannica\_illustrated}
\item \textbf{Dataset Revision:} \texttt{309f765b1ecc83e0a7f7762e44b4a927a7555443}
\item \textbf{Page ID:} \texttt{britannica:309f765b1ecc83e0a7f7762e44b4a927a7555443:train:2077} \\
  \textbf{Filename:} \texttt{page-2077-b2d54e4a68a4.jpg} \\
  \textbf{Row Index:} 2077 \\
  \textbf{SHA256:} \texttt{b2d54e4a68a40bc5...}
\end{itemize}

\begin{figure}[h]
\centering
\includegraphics[width=0.85\textwidth]{../scans/page-2077-b2d54e4a68a4.jpg}
\caption*{Source scan: \texttt{page-2077-b2d54e4a68a4.jpg}}
\end{figure}

\end{document}
